\title{Quantum Monte Carlo analysis of role of interlayer coupling in an organic squaric acid crystal}
\author{Vikas Vijigiri}
\email{vikasvikki@iopb.res.in}
\affiliation{Institute of Physics, Bhubaneswar-751005, Orissa, India}
\affiliation{Homi Bhabha National Institute, Mumbai - 400 094, Maharashtra, India}
\author{Ashis K. Nandy}\email{aknandy@niser.ac.in}\affiliation{School of Physical Sciences, National Institute of Science Education and Research, Jatni, 752050, Odisha, India}
\affiliation{Homi Bhabha National Institute, Mumbai - 400 094, Maharashtra, India}
\author{Saptarshi Mandal }\email{saptarshi@iopb.res.in}\affiliation{Institute of Physics, Bhubaneswar-751005, Orissa, India}
\affiliation{Homi Bhabha National Institute, Mumbai - 400 094, Maharashtra, India}


\begin{abstract}
%We explore the zero-temperature phase diagram of the three-dimensional squaric acid crystal using discrete-time(Suzuki-Trotter decomposition) path integral quantum Monte Carlo technique. We consider a model Hamiltonian with the leading term consisting of a four-spin interacting term, $J_0$, within the plaquette, which along with the Intramolecular coupling term, $J_1$, conserves the ice-rules. Here, we extend the model to include the interlayer exchange coupling, $J_3$, to suit the realistic nature of the material. This work aims to characterize the phases and their corresponding critical boundaries using pseudo-spin formalism.
%We perform our calculations for a generic interlayer coupling with the nature of the interaction is taken at random. Our analysis on order parameter($P,\rho$) and the imaginary-time correlation function $c_\tau$ suggests the existence of an intermediate state with two distinct phases co-existing together. One being the ferroelectric global order and the other a deconfined phase with algebraic correlations.  As a special case, we also consider a uniform (anti)ferromagnetic coupling and observe that the material only gets more stabilized, as the phase boundary now starts shifting towards the right. Finally, we conclude our results by drawing the critical boundary with interlayer coupling as a function of the external transverse field $K_x$ both for random and uniform versions. We also give the results of dynamic structure factor $s^{xx}$ calculated within the Linear-spin wave regime for the 2D version and discuss the relevance of it in brief.

We introduce a three-dimensional pseudo-spin model for  squaric acid crystal ($\text{H}_2\text{SQ}$) which is obtained by  considering an Ising type interlayer coupling ($J_3$) between the two-dimensional layers \cite{vijigiri2018classical} of the squaric acid crystal. To explore the ice-rule dominated physics and its implications, we consider three types of interlayer coupling, i.e, ferroelectric, antiferroelectric and disordered. Firstly, within the spin-wave approximation we find that the spectrum and the spin wave velocity does not depend on the sign of interlayer coupling. Instead, the spin-wave velocity is isotropic and proportional to the magnitude of $J_3$. In the quantum regime, we analyze the model using a Suzuki-Trotter quantum monte carlo method and evaluate the imaginary-time correlation function, susceptibility and specific heat in respective parameter space. The phase diagram as a function of temperature (or external field) shows that the intermediate correlated phase with ice-rule dominated physics survives locally and gets extended for higher temperatures for FM interlayer coupling. Where as for AFM or disordered interlayer coupling the intermediate phase is qualitatively similar to individual two dimensional layer. This points out that strong quantum fluctuations for AFM or disordered coupling weakens the effect of interlayer coupling. Moreover, the transition temperature is decreased upon increasing in interlayer strength for AFM and disordered case, vice-versa for FM case. Finally, we discuss possible experimental ramifications of our analysis.
\end{abstract}


\date{\today}

\pacs{
77.80.-e,
75.10.Jm,
75.40.Mg,
77.84.Fa
%71.10.Pm
%03.65.Ud
%03.65.Vf
%03.67.Mn
}
\maketitle
